\section{From a Scale Band to Recovery}\label{sec:recovery_argument}

The corner test identifies variances that preserve the candidate
barrier for a given observation. A recovery guarantee also needs
a single band that works uniformly along the intermediate trajectory.
We verify the geometric tests for Theorem~\ref{thm:ind_noise}.
Supplementary Sections S4 and S5 certify the numerical margins over
the full parameter box.


\begin{figure*}[t]
\centering
\includegraphics[width=.91\textwidth]{figures/scale_selection.pdf}
\caption{Frozen scale selection across the prespecified geometry grid.
(a) Center of each selected width-$.02$ band.
(b) Largest observed first-entry call for the 32 class-1 draws
under the observation preset, followed by the certified upper bound.
Gray marks settings where the candidate search yields no certificate.
All runs use the common 25-call prefix and 106-call total budget.}
\label{fig:selected_scales}
\end{figure*}

\begin{table}[t]
\centering\small
\setlength{\tabcolsep}{3pt}
\begin{tabular}{@{}p{.31\columnwidth}p{.64\columnwidth}@{}}
\toprule
Result & Guarantee\\
\midrule
Prop.~\ref{prop:two_endpoints} & Endpoint replacement before the midpoint for two aligned templates\\
Props.~\ref{prop:mass_band}--\ref{prop:critical_mass} & Fixed-state mass, drift, or endpoint tests\\
Channel tests below & Along-path progress and controlled terminal entry\\
Thm.~\ref{thm:ind_noise} & Unconditional risk separation, with all classes checked for the upper bound\\
\bottomrule
\end{tabular}
\caption{The scale tests answer different questions. One-step
improvement, violation of a candidate barrier, and full recovery
require separate checks.}
\label{tab:test_scope}
\end{table}

\subsection{Selecting candidate bands}
The barrier cutoffs come from upper bounds on the horizontal
posterior mean at the right face. The low-variance bound drops
the competitor; the high-variance bound retains a lower bound
on its odds. In the channel, bound the competitor's responsibility
uniformly by $\bar w$ and the posterior mean derivative by
$[1/4+a(a+2)\bar w]/v_-$. A positive surplus at the entrance
then gives the affine progress comparison of
Theorem~\ref{thm:criterion}. The transverse faces and the first-step
overshoot must be checked, as must preservation of the prepared
strip through the scheduled switch. For every other class, check
entry separately. At the terminal rectangles, minimum squared-distance
gaps occur at corners and control the leakage terms in the
invariance and refinement bounds. Supplement S5 gives the
explicit inequalities used for selection.

The critical-mass threshold in Proposition~\ref{prop:critical_mass}
is a local endpoint comparison. Its nominal mass ceiling is
$.46385$, while the risk theorem uses $h\in[1.60,1.96]$.
We use the full drift and channel tests in that larger-mass regime.
Table~\ref{tab:test_scope} summarizes the implications. In the
earlier 1681-point diagnostic (Supplement S21), 1383 settings pass
the class-1 channel test but only 1176 of those recover all classes.
The other 207 illustrate the need for the remaining tests.

\subsection{Certified instance}
For \eqref{eq:recovery_box}, the barrier uses
$\xi=.3,\eta_B=.01$, restricted cutoffs $.17,2.5$, and the
channel uses $s_-=-.01,s_E=.53,s_U=1.1,\eta_C=.01,t_+=.05$
and $[v_-,v_+]=[.69,.71]$.
The wider event \eqref{eq:ind_event} uses separate coordinate
tolerances, as detailed in Supplement S4.
Uniform bounds $\partial_s m_v\le.56$ and entrance surplus
$\Delta\ge.016$ imply the comparison
\begin{align}
 \underline s_{k+1}&=.89011\underline s_k+
 (1-.89011)\left(.53+\frac{.016-\delta/(\gamma c_-)}{.44}\right),
 \nonumber\\[-2pt]
 \underline s_0&=-.01,\qquad \underline s_{25}>.5341 .
 \label{eq:main_comparison}
\end{align}
The first coordinate at the crossing is below $.649$. Classes 0 and 2 start in their
terminal rectangles. The rectangles
$G_0=[-.1,.4]\times[-.1,.15]$,
$G_1=[.53,1.1]\times[-.1,.15]$ and
$G_2=[-1.12,.01]\times[1.73,2.27]$ have gaps at least $.06$,
giving leakage at most $4e^{-.03/v_L}$.
\subsection{Risk of the actual terminal iterate}
The terminal center is
$x_{\rm aff}=(B+c_LI)^{-1}(b+c_L\mu_Z)$.
The preceding tests give on $E$ a residual at most
$6q^n+12e^{-.03/v_L}+4\sqrt2\delta_L/c_-$ with $q=.751$.
Its reference risk is
$\sum_{\beta\in\{\alpha,1\}}(\beta\nu^2+c_L^2\tau^2)/(\beta+c_L)^2$.
A budget-uniform fourth-moment envelope controls $E^c$ by
Cauchy--Schwarz. Restricted iterates instead remain at directional
distance at least $.7$ from $\mu_1$ on $E$.
These estimates establish the risk interface; full tail and
residual calculations are in Supplement S5.
